% \cventry{role}{institution}{dates} sets an entry header line, with the dates
% right-aligned in an 8em box; the cvbody environment holds its description.
% Defined here with \providecommand / \ifdefined rather than in
% Includes/Settings.tex so that a document reusing these section files with its
% own preamble also gets them, from whichever file it reads first.
\providecommand{\cventrysep}{0.95em}
\providecommand{\cvfirstentry}{1}
\providecommand{\cventry}[3]{%
  \ifnum\cvfirstentry=0 \par\addvspace{\cventrysep}\fi
  \gdef\cvfirstentry{0}%
  \par\noindent
  \parbox[t]{\dimexpr\linewidth-8.5em\relax}{\raggedright\textbf{#1}, \textit{#2}}%
  \hfill\makebox[8em][r]{\footnotesize\color{black!60}#3}%
  \par\nobreak\vspace{0.1em}}
\ifdefined\cvbody\else
  \newenvironment{cvbody}
    {\begin{list}{}{\setlength{\leftmargin}{1.2em}\setlength{\rightmargin}{0mm}%
       \setlength{\topsep}{0.1em}\setlength{\parsep}{0pt}\setlength{\itemsep}{0pt}}%
     \item[]}
    {\end{list}}
\fi

\section{Current Position}
\gdef\cvfirstentry{1}

\cventry{Associate Professor}{Linköping University}{Aug. 2025 -- present}
\begin{cvbody}
Division of Applied Thermodynamics and Fluid Mechanics, Department of Management and Engineering. My research combines data-driven methods and artificial intelligence with Computational Fluid Dynamics (CFD) to develop efficient and reliable methods for simulation, model reduction, control, and uncertainty quantification of fluid flows. I received my Docent title in Fluid Mechanics at LiU in February 2026.
\end{cvbody}

\section{Previous Appointments}
\gdef\cvfirstentry{1}

\cventry{Researcher}{\href{https://chalmersindustriteknik.se/en/}{Chalmers Industriteknik}, based at Chalmers University of Technology}{2023 -- 2025}
\begin{cvbody}
Employed by Chalmers Industriteknik (CIT), with the research carried out full-time at Chalmers University of Technology. The work applied artificial intelligence and machine learning to computational fluid dynamics, within the project ``Artificial Intelligence in Fluid Dynamics for Improving the Lifetime of Hydraulic Turbines''.
\end{cvbody}

\cventry{Postdoctoral Researcher}{Chalmers University of Technology}{2019 -- 2023}
\begin{cvbody}
Project: New operating procedures of hydropower for a sustainable energy system. Developed advanced numerical methods in OpenFOAM to simulate the transient operation of hydraulic turbines, and expanded my expertise in data-driven and machine-learning approaches, particularly reduced-order modeling such as Proper Orthogonal Decomposition (POD) and Dynamic Mode Decomposition (DMD). The work produced several publications in high-impact peer-reviewed journals and an open-source package for hydraulic turbine transients, now used by Vattenfall.
\end{cvbody}

\cventry{Lecturer}{University of Science and Culture, Tehran, Iran}{2018 -- 2019}
\begin{cvbody}
Taught ``Introduction to Turbomachinery'' at the undergraduate level.
\end{cvbody}
